\chapter{Algebry tensorowa, symetryczna i zewnętrzna}
\label{ch:lin10-algebry}
\index{algebra tensorowa@Algebra tensorowa}
\index{algebra symetryczna@Algebra symetryczna}
\index{algebra zewnętrzna@Algebra zewnętrzna}
\index{iloczyn zewnętrzny@Iloczyn zewnętrzny}

\noindent\textbf{Cel rozdziału.}
Iloczyn tensorowy z Rozdziału~\ref{ch:lin9-iloczyny-tensorowe}
umożliwia linearyzację odwzorowań wieloliniowych.
W tym rozdziale łączymy wszystkie potęgi tensorowe
w jedną algebrę. Dwa ilorazy tej algebry pozwalają
wymusić przemienność czynników albo znikanie iloczynu
dwóch jednakowych wektorów. W ten sposób powstaną
algebry symetryczna i zewnętrzna. Wyprowadzimy ich
własności uniwersalne, bazy i wymiary, a także
związek algebry zewnętrznej z wyznacznikiem
i minorami macierzy.

Przez cały rozdział \(K\) oznacza ciało, \(0_K\ne1_K\).
Nie zakładamy, że \(\operatorname{char}K=0\).
Własności konstrukcyjne dotyczą dowolnych przestrzeni
wektorowych nad \(K\); przy wyliczaniu wymiarów
wyraźnie zakładamy skończony wymiar. Przyjmujemy
\(V^{\otimes0}=K\) oraz \(V^{\otimes1}=V\), a w potęgach
wyższych ustalamy lewostronne nawiasowanie.
Kanoniczne izomorfizmy łączności z
Twierdzenia~\ref{thm:lin9-lacznosc} umożliwiają
przechodzenie między różnymi nawiasowaniami.

\section{Algebry stopniowane i potęgi tensorowe}
\label{sec:lin10-podstawy}

\begin{definicja}[Algebra, ideał i stopniowanie]
\label{def:lin10-algebra}
\emph{Łączną algebrą z jedynką nad \(K\)}
nazywamy przestrzeń wektorową \(A\) nad \(K\)
z dwuliniowym mnożeniem
\(\mu\colon A\times A\to A\), dla którego
\((ab)c=a(bc)\), oraz z elementem \(1_A\)
spełniającym \(1_Aa=a1_A=a\).
Nie wymagamy przemienności mnożenia.
\emph{Homomorfizm algebr z jedynką}
jest odwzorowaniem liniowym zachowującym
mnożenie i element jednostkowy.

Podprzestrzeń \(I\subseteq A\)
nazywamy \emph{ideałem obustronnym},
jeśli \(aI\subseteq I\) i \(Ia\subseteq I\)
dla każdego \(a\in A\).
Algebra \(A\) jest \emph{nieujemnie stopniowana},
jeśli
\(A=\bigoplus_{k\geq0} A^k\) jako przestrzeń
wektorowa i \(A^pA^q\subseteq A^{p+q}\).
W sumie prostej dopuszczamy tylko elementy
o skończenie wielu niezerowych składowych.
\end{definicja}

\begin{stwierdzenie}[Algebra ilorazowa]
\label{prop:lin10-iloraz-algebry}
Niech \(A\) będzie łączną algebrą z jedynką,
a \(I\subseteq A\) ideałem obustronnym.
Na przestrzeni \(A/I\) istnieje dokładnie jedno
mnożenie spełniające
\[
 (a+I)(b+I)=ab+I\qquad(a,b\in A).
\]
Jest ono dwuliniowe i łączne; jedynką jest
\(1_A+I\). Jeśli \(I=A\), iloraz jest
algebrą zerową, której jedynka jest zerem.
Jeżeli \(A\) jest stopniowana, a
\(I=\bigoplus_{k\geq0}(I\cap A^k)\), to
\[
 A/I\cong\bigoplus_{k\geq0}
       A^k/(I\cap A^k)
\]
jako algebra stopniowana.
\end{stwierdzenie}
\begin{proof}
Jeśli \(a'=a+i\), \(b'=b+j\)
dla \(i,j\in I\), to
\(a'b'-ab=aj+ib+ij\in I\).
Mnożenie klas jest więc niezależne od
reprezentantów. Dwuliniowość i łączność
przechodzą z \(A\) do ilorazu;
\(1_A+I\) spełnia obie tożsamości jednostki.
Jednoznaczność wynika ze wzoru na iloczyn klas.

W przypadku stopniowanym każdą klasę
reprezentuje skończona suma
jednorodnych składników. Jeśli
\(\sum_k a_k\in I\), jednorodność ideału
wymusza \(a_k\in I\cap A^k\)
dla każdego \(k\). To dowodzi zarówno
rozpiętości, jak i jednoznaczności
rozkładu klas na składowe stopni.
Mnożenie respektuje stopnie,
ponieważ \(A^pA^q\subseteq A^{p+q}\).
\end{proof}

\begin{definicja}[Algebra tensorowa]
\label{def:lin10-tensorowa}
Dla przestrzeni wektorowej \(V\) nad \(K\)
\emph{algebrą tensorową} nazywamy
stopniowaną przestrzeń
\[
 T(V):=\bigoplus_{k=0}^{\infty}V^{\otimes k}.
\]
Mnożenie na składowych jednorodnych
\(x\in V^{\otimes p}\),
\(y\in V^{\otimes q}\) definiujemy
przez kanoniczne odwzorowanie
\[
 V^{\otimes p}\otimes_K V^{\otimes q}
 \longrightarrow V^{\otimes(p+q)},
 \qquad x\otimes y\longmapsto xy,
\]
odpowiadające złączeniu czynników.
W szczególności
\[
 (v_1\otimes\cdots\otimes v_p)
 (w_1\otimes\cdots\otimes w_q)
 =v_1\otimes\cdots\otimes v_p
  \otimes w_1\otimes\cdots\otimes w_q.
\]
Mnożenie przedłużamy dwuliniowo na
skończone sumy składowych.
Elementem jednostkowym jest \(1_K\in V^{\otimes0}\).
\end{definicja}

\begin{twierdzenie}[Własność uniwersalna algebry tensorowej]
\label{thm:lin10-tensorowa-uniwersalna}
Dla dowolnej przestrzeni \(V\)
algebra \(T(V)\) jest łączną algebrą z jedynką.
Dla każdej łącznej algebry \(A\) z jedynką
i każdego odwzorowania liniowego
\(f\colon V\to A\)
istnieje dokładnie jeden homomorfizm
algebr z jedynką
\(\widehat f\colon T(V)\to A\) taki, że
\(\widehat f(v)=f(v)\) dla każdego \(v\in V\).
Na tensorach prostych ma on postać
\[
 \widehat f(v_1\otimes\cdots\otimes v_k)
 =f(v_1)\cdots f(v_k),\qquad
 \widehat f(a)=a1_A\quad(a\in V^{\otimes0}=K).
\]
\end{twierdzenie}
\begin{proof}
Mnożenie jednorodnych tensorów jest dobrze
określone dzięki własności uniwersalnej
iloczynu tensorowego i kanonicznym
izomorfizmom łączności
(Twierdzenie~\ref{thm:lin9-lacznosc}).
Na tensorach prostych oba nawiasowania
\((xy)z\) i \(x(yz)\) zapisują czynniki
w identycznej kolejności.
Ponieważ tensory proste rozpinają
każdą potęgę tensorową, równość rozszerza
się na wszystkie składowe jednorodne,
a następnie przez dwuliniowość na \(T(V)\).
Analogicznie konkatenacja z tensorem
pustym daje \(1_Kx=x1_K=x\).
Skończoność liczby niezerowych składowych
gwarantuje, że każdy otrzymany iloczyn
należy do algebraicznej sumy prostej.

Ustalmy \(f\colon V\to A\).
Dla \(k\geq1\) odwzorowanie
\((v_1,\ldots,v_k)\mapsto f(v_1)\cdots f(v_k)\)
jest \(k\)-liniowe: przy ustaleniu
pozostałych argumentów jest złożeniem
odwzorowania liniowego \(f\) i
mnożenia przez ustalone elementy \(A\).
Twierdzenie~\ref{thm:lin9-wielokrotny}
daje jedyne liniowe
\(f_k\colon V^{\otimes k}\to A\)
o wskazanych wartościach.
W stopniu \(0\) określamy
\(f_0(a)=a1_A\).
Dla \(t=\sum_{k=0}^{N}t_k\)
definiujemy
\(\widehat f(t)=\sum_{k=0}^{N}f_k(t_k)\).
Rozkład na składowe jest jednoznaczny,
więc definicja jest dobra i liniowa.
Na dowolnych tensorach prostych mamy
\(\widehat f(xy)=\widehat f(x)\widehat f(y)\);
dwuliniowość przenosi tę równość
na wszystkie pary \(x,y\in T(V)\).
Ponadto \(\widehat f(1_K)=1_A\),
więc \(\widehat f\) jest homomorfizmem
z jedynką.

Jeśli \(g\colon T(V)\to A\) jest drugim
takim homomorfizmem,
to \(g(v_1\otimes\cdots\otimes v_k)
=g(v_1)\cdots g(v_k)
=f(v_1)\cdots f(v_k)\).
W stopniu \(0\) każdy homomorfizm
z jedynką spełnia \(g(a)=a1_A\).
Ponieważ tensory proste i składowa
\(K\) rozpinają \(T(V)\),
otrzymujemy \(g=\widehat f\).
\end{proof}

\begin{przyklad}[Dlaczego algebra tensorowa nie jest przemienna]
\label{prz:lin10-nieprzemienna}
Niech \(V=\mathbb R^2\) z bazą
\((e_1,e_2)\), a \(A=\operatorname{Mat}_{2\times2}(\mathbb R)\).
Zadajmy \(f(e_1)=E_{12}\),
\(f(e_2)=E_{21}\), gdzie \(E_{ij}\)
są jednostkami macierzowymi.
Własność uniwersalna daje homomorfizm
\(\widehat f\colon T(V)\to A\), zatem
\[
 \widehat f(e_1e_2)=E_{11},\qquad
 \widehat f(e_2e_1)=E_{22}.
\]
Wyniki są różne, więc
\(e_1e_2\ne e_2e_1\) w \(T(V)\).
Porządek czynników zachowuje informację.
Przemienność można dopiero wymusić,
przechodząc do odpowiedniego ilorazu.
\end{przyklad}

\section{Algebra symetryczna}
\label{sec:lin10-symetryczna}

\begin{definicja}[Algebra symetryczna]
\label{def:lin10-symetryczna}
Niech \(J\) będzie ideałem obustronnym
\(T(V)\) generowanym przez wszystkie
\(u\otimes v-v\otimes u\), gdzie \(u,v\in V\).
\emph{Algebrą symetryczną} nazywamy
\[
 S(V):=T(V)/J.
\]
Iloczyn klas oznaczamy zwykłym mnożeniem.
Składową stopnia \(k\) oznaczamy
\(S^k(V):=V^{\otimes k}/(J\cap V^{\otimes k})\).
W szczególności \(S^0(V)=K\),
\(S^1(V)=V\), a \(S(V)\) jest
algebrą przemienną.
\end{definicja}

\begin{stwierdzenie}[Stopniowanie i przemienność]
\label{prop:lin10-sym-stopnie}
Ideał \(J\) jest jednorodny, więc
\(S(V)=\bigoplus_{k\geq0}S^k(V)\).
Dla wszystkich \(u,v\in V\) w ilorazie
zachodzi \(uv=vu\); mnożenie w \(S(V)\)
jest przemienne, łączne i posiada jedynkę.
\end{stwierdzenie}
\begin{proof}
Każdy generator \(u\otimes v-v\otimes u\)
leży w \(V^{\otimes2}\).
Dowolny element ideału \(J\)
jest skończoną sumą elementów
\(a(u\otimes v-v\otimes u)b\).
Rozkładając \(a,b\) na jednorodne
składowe, otrzymujemy składniki
jednorodne należące do \(J\).
Stąd \(J=\bigoplus_{k\ge0}(J\cap V^{\otimes k})\).
Twierdzenie~\ref{prop:lin10-iloraz-algebry}
daje stopniowanie i łączność.
W ilorazie obrazy generatorów
\(u,v\in V\) komutują.
Każdy tensor prosty jest iloczynem
takich generatorów; zamiany sąsiednich
czynników pozwalają przestawić
dowolne dwa jednomiany. Z liniowości
wynika przemienność wszystkich
elementów \(S(V)\).
Ponieważ \(J\) ma składniki
wyłącznie stopnia co najmniej \(2\),
klasy w stopniach \(0,1\) nie są
dodatkowo utożsamiane.
\end{proof}

\begin{twierdzenie}[Własność uniwersalna algebry symetrycznej]
\label{thm:lin10-sym-uniwersalna}
Niech \(A\) będzie przemienną,
łączną algebrą z jedynką nad \(K\)
i \(f\colon V\to A\) odwzorowaniem
liniowym. Istnieje dokładnie jeden
homomorfizm algebr z jedynką
\(\overline f\colon S(V)\to A\)
taki, że \(\overline f(v)=f(v)\)
dla wszystkich \(v\in S^1(V)=V\).
\end{twierdzenie}
\begin{proof}
Z Twierdzenia~\ref{thm:lin10-tensorowa-uniwersalna}
otrzymujemy jedyny homomorfizm
\(\widehat f\colon T(V)\to A\)
przedłużający \(f\).
Ponieważ mnożenie w \(A\)
jest przemienne,
\[
 \widehat f(u\otimes v-v\otimes u)
 =f(u)f(v)-f(v)f(u)=0.
\]
Jądro homomorfizmu jest
ideałem obustronnym, a więc
zawiera ideał \(J\) generowany
przez te różnice.
Wobec tego wzór
\(\overline f(t+J)=\widehat f(t)\)
jest dobrze określony, liniowy
i zachowuje mnożenie oraz jedynkę.
Przedłuża \(f\).
Jeśli \(g\) jest innym takim
homomorfizmem, to
\(g(v_1\cdots v_k)
=f(v_1)\cdots f(v_k)\)
i \(g(a)=a1_A\) dla \(a\in K\).
Ponieważ takie iloczyny
rozpinają \(S(V)\), mamy
\(g=\overline f\).
\end{proof}

\begin{twierdzenie}[Baza potęg symetrycznych]
\label{thm:lin10-sym-baza}
Niech \(V\) będzie skończenie
wymiarową przestrzenią nad \(K\),
\(n=\dim V\), z uporządkowaną
bazą \(B=(e_1,\ldots,e_n)\).
Dla każdego \(k\ge0\)
rodzina jednomianów
\[
 e_1^{a_1}\cdots e_n^{a_n},
 \qquad a_i\in\mathbb N_0,\quad
 \sum_{i=1}^n a_i=k,
\]
jest bazą przestrzeni \(S^k(V)\).
Gdy \(n\ge1\), zachodzi
\[
 \dim S^k(V)=\binom{n+k-1}{k}.
\]
Gdy \(n=0\), przestrzeń \(S^0(V)=K\),
a \(S^k(V)=\{0\}\) dla \(k\ge1\).
Wszystkie te twierdzenia obowiązują
w dowolnej charakterystyce ciała.
\end{twierdzenie}
\begin{proof}
Każdy tensor prosty jest iloczynem
wektorów \(V\), a każdy z nich
jest kombinacją liniową \(e_i\).
Po rozwinięciu przez wieloliniowość
i przestawieniu czynników w \(S(V)\)
otrzymujemy kombinację wskazanych
jednomianów; rozpinają one \(S^k(V)\).

Dla niezależności zbudujmy
homomorfizm
\(P\colon S(V)\to K[x_1,\ldots,x_n]\).
Niech \(f(e_i)=x_i\);
jest to dobrze określone
odwzorowanie liniowe \(V\to K[x_1,\ldots,x_n]\).
Własność uniwersalna
z Twierdzenia~\ref{thm:lin10-sym-uniwersalna}
daje homomorfizm \(P\)
spełniający
\(P(e_1^{a_1}\cdots e_n^{a_n})
=x_1^{a_1}\cdots x_n^{a_n}\).
Z definicji pierścienia
wielomianów różne jednomiany
są liniowo niezależne nad \(K\).
Zatem jeśli kombinacja wskazanych
wektorów w \(S^k(V)\) jest zerowa,
to wszystkie jej współczynniki
są zerowe po zastosowaniu \(P\).
To dowodzi niezależności.

Dla \(n\ge1\) liczba
\(n\)-tek nieujemnych liczb całkowitych
o sumie \(k\) jest
\(\binom{n+k-1}{k}\):
można je zakodować,
umieszczając \(n-1\) przegród
wśród \(n+k-1\) miejsc.
Nie używaliśmy dzielenia
przez żaden element ciała,
więc twierdzenie obejmuje
również dodatnią charakterystykę.
Przy \(n=0\) jedyna krotka
wykładników jest pusta
i ma sumę \(0\); daje
bazę \(S^0(V)=K\).
Dla \(k>0\) nie ma
żadnego jednomianu,
co zgodnie z rozpiętością
oznacza przestrzeń zerową.
\end{proof}

\begin{wniosek}[Model wielomianowy]
\label{cor:lin10-wielomiany}
Przy założeniach
Twierdzenia~\ref{thm:lin10-sym-baza}
istnieje izomorfizm algebr
\[
 S(V)\cong K[x_1,\ldots,x_n],
 \qquad e_i\longmapsto x_i,
\]
zachowujący stopnie.
Dla \(n=0\) prawa strona
oznacza \(K\).
\end{wniosek}
\begin{proof}
Homomorfizm \(P\) skonstruowany
w poprzednim dowodzie
przekształca bazę każdej
składowej stopnia \(k\)
bijektywnie na bazę jednomianów
jednorodnych stopnia \(k\).
Jest więc bijekcją liniową
na każdej składowej, a zatem
na algebraicznej sumie prostej.
Jako homomorfizm z jedynką
jest izomorfizmem algebr.
\end{proof}

\begin{przyklad}[Potęga symetryczna przestrzeni dwuwymiarowej]
\label{prz:lin10-sym2}
Niech \(V=K^2\) z bazą \((e_1,e_2)\).
Wówczas
\(S^2(V)=\operatorname{span}
 \{e_1^2,e_1e_2,e_2^2\}\)
i \(\dim S^2(V)=3\).
W \(S(V)\) mamy \(e_1e_2=e_2e_1\),
podczas gdy w \(V\otimes V\)
oba tensory są różne.
W charakterystyce \(2\)
zachodzi \((e_1+e_2)^2=e_1^2+e_2^2\);
nie oznacza to zaniku
elementu \(e_1e_2\), który
nadal należy do wskazanej bazy.
\end{przyklad}
